% Class-compatibility preamble retained for a reader excerpt of
%   Dave Witte Morris, "A Cohomological Proof that Real Representations of
%   Semisimple Lie Algebras Have Q-Forms", SIGMA 11 (2015) 034,
%   DOI 10.3842/SIGMA.2015.034, CC BY 4.0.
%
% The source pins the publisher class sigma, which is not part of the delivery
% baseline, so this file supplies the amsart baseline and then copies the author's
% own declaration block mechanically and unchanged:
%   - pinned-original-source.tex lines 11-31: the author's DeclareMathOperator
%     list, the equation numbering and the author's newtheorem declarations;
%   - pinned-original-source.tex lines 35-43, which the author writes immediately
%     after \begin{document} in the publisher layout: the author's own \Res,
%     \ResOne and \Lie macros and \allowdisplaybreaks.
% The publisher header commands of the sigma class (\FirstPageHeading,
% \ArticleName and the like) and the macro \PaperNumber they require are not
% reproduced, because this excerpt carries neither the journal front matter nor
% anything else that uses them.
% No mathematical macro, statement, symbol or assumption is introduced, changed
% or removed by this file.
\documentclass[10pt,a4paper]{amsart}
\usepackage{amssymb,amsmath,amsthm}
\usepackage{enumerate}
\usepackage{hyperref}

\DeclareMathOperator{\Cliff}{Clif\/f}
\DeclareMathOperator{\Gal}{Gal}
\DeclareMathOperator{\Inn}{Inn}
\DeclareMathOperator{\Mat}{Mat}
\DeclareMathOperator{\Out}{Out}
\numberwithin{equation}{section}

\newtheorem{prop}{Proposition}[section]
\newtheorem{cor}[prop]{Corollary}
\newtheorem{lem}[prop]{Lemma}
{\theoremstyle{definition}
\newtheorem{defn}[prop]{Definition}
\newtheorem{notation}[prop]{Notation}
\newtheorem{rem}[prop]{Remark}
\newtheorem{ques}[prop]{Question}
\newtheorem{eg}[prop]{Example}
\newtheorem{erratum}[prop]{Erratum}}

\newcommand{\Res}[1]{\mathop{\mathbf{R}_{#1}}}
\newcommand{\ResOne}[1]{\mathop{\mathbf{R}^{(1)}_{#1}}}
\newcommand{\Lie}[1]{\mathfrak{\lowercase{#1}}}
\allowdisplaybreaks
\title{Four equivalent descriptions of R-universality of a semisimple Lie algebra over Q}
\author{Dave Witte Morris}
\date{doi:10.3842/SIGMA.2015.034; CC BY 4.0}
\begin{document}
\maketitle
\noindent\textit{Editorial note.} All delivered mathematical text is the author's own, copied from the pinned publisher source sources/algebra-fl02-a-cohomological-proof-that-real-re/pinned-original-source.tex; no statement or proof step was rewritten, reordered or removed.\par\medskip
\noindent\textit{Editorial note.} Preamble provenance: the source pins the publisher class sigma, which is not part of the delivery baseline, so this card uses sources/algebra-fl02-a-cohomological-proof-that-real-re/retained-preamble.tex, which supplies the amsart baseline, copies the author's DeclareMathOperator list, equation numbering and newtheorem declarations from pinned source lines 11-31, and copies the author's own \textbackslash{}Res, \textbackslash{}ResOne and \textbackslash{}Lie macros from pinned source lines 35-43, which the author writes immediately after \textbackslash{}begin\{document\}. No macro was invented and no mathematical symbol was changed.\par\medskip
\noindent\textit{Editorial note.} Bibliography provenance: the counted statement and its proof contain no citation, so this card carries no bibliography. The source has a thebibliography environment at pinned source lines 954-1016 whose three entries are used elsewhere in the paper; nothing from it is needed here and no entry was reconstructed.\par\medskip
\noindent\textit{Editorial note.} External references are navigation pointers only; every deliberately kept step of the author's argument stays inside the counted spans. Their printed numbers are the author's own numbering under the author's newtheorem declarations, computed mechanically by the lane's scanner.\par\medskip
\noindent\textit{Allowed standard prerequisites (editorial).}\begin{itemize}
\item The Tits algebras D(Lie G, K)(lambda) attached to a dominant weight and the representations rho\_lambda over Q and over R. Set up in the author's introduction and in the definitions preceding the counted statement (source lines 79-201). Supplies the four objects that the four equivalent conditions speak about.
\end{itemize}\medskip
\begin{lem}
\label{RunivIff}
If $\Lie G$ is a~semisimple Lie algebra over~$\mathbb{Q}$, then the following are equivalent:
\renewcommand{\labelenumi}{{\rm \theenumi.}}
\begin{enumerate}\itemsep=0pt
\item
\label{RunivIff-univ}
$\Lie G$ is $\mathbb{R}$-universal.
\item
\label{RunivIff-irred}
${}^\mathbb{Q} \rho_\lambda \otimes_\mathbb{Q} \mathbb{R} \cong {}^\mathbb{R} \rho_\lambda$, for every dominant
weight~$\lambda$.
\item
\label{RunivIff-DivAlg}
$D_{\Lie G, \mathbb{Q}}(\lambda) \otimes_\mathbb{Q} \mathbb{R}$ is a~division algebra, for every dominant
weight~$\lambda$.
\item
\label{RunivIff-D}
$D_{\Lie G, \mathbb{Q}}(\lambda) \otimes_\mathbb{Q} \mathbb{R} \cong D_{\Lie G, \mathbb{R}}(\lambda)$, for every
dominant weight~$\lambda$.
\end{enumerate}
\end{lem}

\begin{proof}
($\ref{RunivIff-univ} \Rightarrow \ref{RunivIff-irred}$) Let~$\rho$ be a~$\mathbb{Q}$-form of ${}^\mathbb{R}
\rho_\lambda$.
Then~$\rho$ must be $\mathbb{Q}$-irreducible (since $\rho \otimes_\mathbb{Q} \mathbb{R} \cong {}^\mathbb{R}
\rho_\lambda$ is $\mathbb{R}$-irreducible) and
\begin{gather*}
\rho \otimes_\mathbb{Q} \mathbb{C} \cong (\rho \otimes_\mathbb{Q} \mathbb{R}) \otimes_\mathbb{R} \mathbb{C} \cong
{}^\mathbb{R} \rho_\lambda \otimes_\mathbb{R} \mathbb{C}
\end{gather*}
has an irreducible summand with highest weight~$\lambda$.
Therefore $\rho \cong {}^\mathbb{Q} \rho_\lambda$.

($\ref{RunivIff-irred} \Rightarrow \ref{RunivIff-univ}$) ${}^\mathbb{Q} \rho_\lambda$ is a~$\mathbb{Q}$-form of
${}^\mathbb{R} \rho_\lambda$.

($\ref{RunivIff-irred} \Leftrightarrow \ref{RunivIff-DivAlg} \Leftrightarrow \ref{RunivIff-D}$) We have
$\mathrm{End}_{\Lie G}({}^\mathbb{Q} \rho_\lambda \otimes_\mathbb{Q} \mathbb{R}) = D_{\Lie G, \mathbb{Q}}(\lambda)
\otimes_\mathbb{Q} \mathbb{R}$.
So ${}^\mathbb{Q} \rho_\lambda \otimes_\mathbb{Q} \mathbb{R}$ is irreducible (and hence equal to~${}^\mathbb{R}
\rho_\lambda$) if and only if $D_{\Lie G, \mathbb{Q}}(\lambda) \otimes_\mathbb{Q} \mathbb{R}$ is a~division algebra (and
hence equal to~$D_{\Lie G, \mathbb{R}}(\lambda)$).
\end{proof}

% Bibliography excerpt retained for reader excerpts of
%   Dave Witte Morris, "A Cohomological Proof that Real Representations of
%   Semisimple Lie Algebras Have Q-Forms", SIGMA 11 (2015) 034,
%   DOI 10.3842/SIGMA.2015.034, CC BY 4.0.
%
% Every entry below is copied character for character from the author's own
% thebibliography environment of the pinned publisher source
% (pinned-original-source.tex lines 954-1016).  No entry was rewritten,
% reordered, merged, split or invented, and no citation key was changed.
\begin{thebibliography}{99}

\bibitem{Morris-RRepsQForms}
Morris D., Real representations of semisimple {L}ie algebras have {${\mathbb
  Q}$}-forms, in Algebraic Groups and Arithmetic, Tata Institute of Fundamental
  Research, Mumbai, 2004, 469--490.


\bibitem{PrasadRapinchuk-Prescribed}
Prasad G., Rapinchuk A., On the existence of isotropic forms of semi-simple
  algebraic groups over number f\/ields with prescribed local behavior,
  \href{http://dx.doi.org/10.1016/j.aim.2006.01.001}{\textit{Adv. Math.}} \textbf{207} (2006), 646--660.


\bibitem{Tits-RepLinIrred}
Tits J., Repr\'esentations lin\'eaires irr\'eductibles d'un groupe r\'eductif
  sur un corps quelconque, \href{http://dx.doi.org/10.1515/crll.1971.247.196}{\textit{J.~Reine Angew. Math.}} \textbf{247} (1971),
  196--220.



\end{thebibliography}
\end{document}
